\documentclass[11pt,reqno]{amsart}
\usepackage[localtheorems]{raeez-math-template}

\usepackage{array}

\newtheorem{theorem}{Theorem}[section]
\newtheorem{proposition}[theorem]{Proposition}
\newtheorem{lemma}[theorem]{Lemma}
\newtheorem{corollary}[theorem]{Corollary}
\theoremstyle{definition}
\newtheorem{definition}[theorem]{Definition}
\newtheorem{construction}[theorem]{Construction}
\newtheorem{example}[theorem]{Example}
\theoremstyle{remark}
\newtheorem{remark}[theorem]{Remark}

\providecommand{\cA}{\mathcal{A}}
\providecommand{\cM}{\mathcal{M}}
\providecommand{\barB}{\bar{B}}
\providecommand{\Defcyc}{\mathrm{Def}_{\mathrm{cyc}}}
\providecommand{\MC}{\mathrm{MC}}
\providecommand{\Hom}{\mathrm{Hom}}
\providecommand{\Ext}{\mathrm{Ext}}
\providecommand{\CE}{\mathrm{CE}}
\providecommand{\Aut}{\mathrm{Aut}}
\providecommand{\fg}{\mathfrak{g}}
\providecommand{\gmod}{\mathfrak{g}^{\mathrm{mod}}_{\cA}}
\providecommand{\gmodzero}{\mathfrak{g}^{\mathrm{mod},(0)}_{\cA}}
\providecommand{\Sh}{\mathrm{Sh}}
\providecommand{\Vir}{\mathrm{Vir}}
\providecommand{\Heis}{\mathrm{Heis}}
\providecommand{\Ash}{\cA^{\mathrm{sh}}}
\providecommand{\Linf}{L_\infty}
\providecommand{\Ainf}{A_\infty}
\providecommand{\Trees}{\mathsf{Trees}}
\providecommand{\cW}{\mathcal{W}}

\numberwithin{equation}{section}

% ================================================================
\begin{document}

\title[Genus-zero shadow and formality obstruction]
{The shadow obstruction tower of a chiral algebra\\
and the transferred $L_\infty$ tower}

\author{Raeez Lorgat}
\address{Department of Mathematics, Harvard University,
 Cambridge, MA 02138, USA}
\email{rlorgat@math.harvard.edu}

\date{}

\begin{abstract}
Let $\cA$ be a chiral algebra on a smooth projective curve and
let $\gmod$ be its modular convolution dg Lie algebra, with
bar-intrinsic Maurer--Cartan element
$\Theta_{\cA}\in\MC(\gmod)$. In the verified chirally Koszul
genus-zero cyclic setting, after a cyclic contraction of
$\gmodzero$ onto its shadow cohomology is fixed, the finite-order
shadow projections of $\Theta_{\cA}$ coincide with the transferred
$\Linf$ bracket evaluations on the projected Maurer--Cartan
element. Degree two is the curvature input; degrees $r\ge 3$ are
the formality obstruction evaluations. For the standard families
where the transfer data are constructed, this comparison identifies
the shadow termination degree with the $\Linf$ obstruction level.
\end{abstract}

\subjclass[2020]{17B69, 18M70, 18N40, 81T40}

\keywords{Chiral algebras, bar construction, $L_\infty$-formality,
homotopy transfer, Maurer--Cartan, operadic complexity}

\maketitle

\section{Introduction}\label{sec:intro}

\subsection{Formality and its obstructions}

A dg Lie algebra $\fg$ is \emph{$\Linf$-formal} if its minimal
$\Linf$-model is $\Linf$-isomorphic to the strict graded Lie algebra
$(H^*(\fg),0,[-,-]_H)$. For a fixed contraction, the transferred
brackets $\ell_r^{\mathrm{tr}}$ are representatives; the obstruction
is their class in the appropriate deformation complex after all
lower obstructions have been killed.
Kontsevich's formality theorem~\cite{Kontsevich03} proved that
polyvector fields on a smooth manifold form an $\Linf$-formal
dg Lie algebra, giving deformation quantization of Poisson
structures; Tamarkin~\cite{Tamarkin03} recovered the result
operadically from formality of the little disks. When formality
fails, obstructions are recorded in cohomology classes
$o_3,o_4,\dots$: at degree $r$ the class $o_r$ measures whether
$\ell_r^{\mathrm{tr}}$ vanishes on cohomology; if $o_r=0$, a
homotopy kills $\ell_r^{\mathrm{tr}}$ and the next obstruction
lives one degree higher. The sequence $(o_r)_{r\ge 3}$ is the
\emph{formality obstruction tower}; see
Loday--Vallette~\cite[Ch.~10]{LodayVallette12}.

\subsection{The shadow tower of a chiral algebra}

On a smooth projective curve $X$, a chiral algebra $\cA$
\cite{BD04} has a bar complex
$\barB(\cA)=T^c(s^{-1}\bar{\cA})$, where $\bar{\cA}=\ker(\epsilon)$
is the augmentation ideal, with deconcatenation coproduct and
differential built from the Arnold form
$\eta_{ij}=d\log(z_i-z_j)$ on the Fulton--MacPherson
compactification of the configuration space. The modular
convolution dg Lie algebra $\gmod$ is the pronilpotent dg Lie
algebra whose differential encodes graph sums over stable
curves in $\overline{\cM}_{g,n}$. The bar-intrinsic element
$\Theta_{\cA}:=D_{\cA}-d_0\in\MC(\gmod)$ satisfies
$D\Theta_{\cA}+\tfrac12[\Theta_{\cA},\Theta_{\cA}]=0$ and
packages the full genus expansion of $\cA$
\cite[MC2 and modular convolution theorem]{Lorgat26}. Its finite-degree truncations
$\Theta^{\le r}_{\cA}$ define the \emph{shadow obstruction
tower}
\[
\Sh_r(\cA)\;:=\;-\nabla_H^{-1}\bigl([\Theta^{\le r-1}_{\cA},
\Theta^{\le r-1}_{\cA}]_r^{(0)}\bigr),
\qquad r\ge 2,
\]
where the bracket is genus-zero graph composition and
$\nabla_H$ is the Hessian covariant derivative on the shadow
algebra $\Ash$. The first three shadows are
$\Sh_2(\cA)=\kappa(\cA)$ (the modular characteristic
\cite[Theorem~D]{Lorgat26}), $\Sh_3(\cA)=\mathfrak{C}(\cA)$
(cubic shadow), and $\Sh_4(\cA)=\mathfrak{Q}(\cA)$ (quartic
shadow).

\subsection{Shadow and transfer}

The two towers are built in different complexes. The formality tower
is the transferred cyclic Chevalley tower of the genus-zero
convolution dg Lie algebra. The shadow tower is the scalar projection
of the bar-intrinsic Maurer--Cartan recursion. The comparison is a
theorem only after the genus-zero cyclic contraction is fixed.

\begin{theorem}[Genus-zero shadow--formality comparison]
\label{thm:n6-shadow-formality-summary}
Let $\cA$ be a chirally Koszul standard algebra for which the
genus-zero cyclic convolution complex, shadow projection, and cyclic
contraction $(\iota,\pi,h)$ are constructed. Assume the cyclic
stable-tree/HPL comparison lemma: after cyclic symmetrisation, the
genus-zero stable-tree graph sum and the Kadeishvili--Merkulov
transfer sum have identical signs, automorphism weights,
propagators, and vertex decorations. For every $r\ge2$,
\[
\Sh_r(\cA)\;=\;\ell_r^{(0),\mathrm{tr}}
\bigl(\pi\Theta^{\le r-1}_{\cA},\ldots,
\pi\Theta^{\le r-1}_{\cA}\bigr)
\quad\text{in }\Ash_{r,0},
\]
relative to the fixed cyclic contraction. Degree $2$ is the
curvature term. For $r\ge3$ this is the transferred formality
obstruction evaluated on the truncated Maurer--Cartan class.
\end{theorem}

The proof has two ingredients. The cases $r=2,3,4$ are the
explicit computations of Proposition~\ref{prop:base-cases}. For
$r\ge 5$, the Kadeishvili--Merkulov tree
formula~\cite{Kadeishvili80,Merkulov99} on the strict genus-zero
convolution algebra projects, after cyclic symmetrisation, to the
same stable-tree graph sum that defines the shadow coefficient.

\subsection{Consequences}

Theorem~\ref{thm:n6-shadow-formality-summary} identifies the shadow termination degree
$r_{\max}(\cA)$ with the $\Linf$ obstruction level
$f_\infty(\cA)$ of $\gmodzero$ for the verified standard cyclic
families. Operation-level equality with a transferred
$\Ainf$-depth requires a separate nondegeneracy statement.
The four shadow-depth classes are realized by Heisenberg (class
\textbf{G}, $r_{\max}=2$, formal), affine Kac--Moody (class
\textbf{L}, $r_{\max}=3$), $\beta\gamma$ (class \textbf{C},
$r_{\max}=4$), and Virasoro (class \textbf{M},
$r_{\max}=\infty$, intrinsically non-formal); see
Section~\ref{sec:examples}.

\subsection*{Scope}

A verified chirally Koszul standard algebra is one for which the
genus-zero cyclic convolution complex, the shadow cohomology
projection, and a cyclic contraction
$(\iota,\pi,h)$ have been constructed. All comparison statements
below are relative to this fixed transfer datum.
They do not assert formality of the full two-coloured
$\mathrm{SC}^{\mathrm{ch,top}}$ operad, nor of the Calabi--Yau
stage-one factorisation algebra before Stage-2 specialisation.

\subsection*{Conventions}

Grading is cohomological with $|d|=+1$; desuspension lowers
degree, $|s^{-1}v|=|v|-1$. Bar complexes use the augmentation
ideal: $\barB(\cA)=T^c(s^{-1}\bar{\cA})$. The propagator
$d\log E(z,w)$ on the bar complex has weight one. All statements
are at the level of cohomology classes unless otherwise noted.

% ================================================================
\section{The shadow obstruction tower}\label{sec:shadow-tower}
% ================================================================

\subsection{The bar complex and its convolution algebra}

Let $X$ be a smooth projective curve over $\mathbb{C}$ and let
$\cA$ be a chiral algebra on $X$~\cite{BD04}. The ordered
chiral bar complex
$\barB(\cA)=T^c(s^{-1}\bar{\cA})
=\bigoplus_{n\ge 0}(s^{-1}\bar{\cA})^{\otimes n}$
carries the deconcatenation coproduct, with differential
assembled from $\eta_{ij}=d\log(z_i-z_j)$ and the chiral product
$\mu\colon j_*j^*(\cA\boxtimes\cA)\to\Delta_*\cA$. The
\emph{modular convolution dg Lie algebra} is
\[
\gmod\;:=\;\Hom^{\mathrm{cyc}}
\bigl(\barB(\cA),\,\cA\bigr)\otimes\mathsf{G},
\]
where $\mathsf{G}=\bigoplus_{g\ge 0}\mathsf{G}^{(g)}$ is the
coefficient algebra generated by stable graphs on
$\overline{\cM}_{g,n}$. The genus filtration gives
$\gmod=\bigoplus_{g\ge 0}(\gmod)_g$; the genus-zero part
$\gmodzero=\Hom^{\mathrm{cyc}}(\barB(\cA),\cA)\otimes
\mathsf{G}^{(0)}$ is a strict dg Lie algebra with bracket
$[-,-]^{(0)}$ from tree composition, since genus-zero stable
graphs are trees.

\subsection{The Maurer--Cartan element and degree truncations}

The differential on $\gmod$ splits as $D=d_0+D_{\mathrm{int}}$,
with $d_0$ linear and $D_{\mathrm{int}}$ the interacting piece
encoding the chiral product and the Arnold form. The difference
$\Theta_{\cA}:=D_{\cA}-d_0\in\gmod$ is a Maurer--Cartan
element~\cite[MC2 bar-intrinsic construction]{Lorgat26}:
\begin{equation}\label{eq:mc-equation}
D\Theta_{\cA}+\tfrac12[\Theta_{\cA},\Theta_{\cA}]=0.
\end{equation}
The element $\Theta_{\cA}$ is bar-intrinsic. Write
\[
F_{\le r}\barB(\cA)=\bigoplus_{n\le r}
(s^{-1}\bar{\cA})^{\otimes n},\qquad
F_{>r}\barB(\cA)=\bigoplus_{n>r}
(s^{-1}\bar{\cA})^{\otimes n}.
\]
The truncation $\Theta^{\le r}_{\cA}$ is the projection of
$\Theta_{\cA}$ modulo $F_{>r}$. The Maurer--Cartan recursion
determines each new degree up to a coboundary from the previous
truncation.

\begin{definition}[Shadow obstruction tower]
\label{def:n6-shadow-obstruction-tower}
The \emph{degree-$r$ shadow} of $\cA$ is
\begin{equation}\label{eq:shadow-definition}
\Sh_r(\cA)\;:=\;-\nabla_H^{-1}\bigl(
[\Theta^{\le r-1}_{\cA},
\Theta^{\le r-1}_{\cA}]_r^{(0)}\bigr)
\;\in\;\Ash_{r,0},
\end{equation}
where the bracket is the genus-zero graph composition, the
subscript $r$ extracts the degree-$r$ component, and $\nabla_H$ is
the Hessian covariant derivative on the shadow algebra
$\Ash=H^*(\gmodzero)$ descended from the convolution Lie bracket.
The sequence $(\Sh_r(\cA))_{r\ge 2}$ is the \emph{shadow
obstruction tower} of $\cA$.
\end{definition}

The definition is bar-intrinsic: if
$\Theta_{\cA}\equiv\Theta'_{\cA}$ mod gauge, then $\Sh_r(\cA)$
and $\Sh_r'(\cA)$ coincide as cohomology classes in $\Ash_{r,0}$.

\begin{remark}[Closed formulas]
On the scalar projection, writing
$\Theta_{\cA}=\kappa\cdot\eta\otimes\Lambda$ and
$H(t)=\sum_{r\ge2}rS_r t^r$, one has
$\Sh_2=\kappa$, $\Sh_3=\mathfrak{C}$, $\Sh_4=\mathfrak{Q}$, and
the shadow metric identity
\[
H(t)^2=t^4\bigl(4\kappa^2+12\kappa\mathfrak{C}\,t
+(9\mathfrak{C}^2+16\kappa\mathfrak{Q})\,t^2\bigr)
\]
on the families where the Riccati reduction has been proved.
\end{remark}

% ================================================================
\section{The $\Linf$-formality obstruction tower}
\label{sec:formality-tower}
% ================================================================

\subsection{Homotopy transfer on a strict model}

Let $\fg$ be a strict dg Lie algebra with differential $d$ and
bracket $[-,-]$. A \emph{deformation retract} onto $H^*(\fg)$ is
a triple $(\iota,\pi,h)$ with
$\iota\colon H^*(\fg)\to\fg$,
$\pi\colon\fg\to H^*(\fg)$, and $h\colon\fg\to\fg$ of degree
$-1$, satisfying $\pi\iota=\mathrm{id}$ and
$\mathrm{id}-\iota\pi=[d,h]$.
The homotopy transfer theorem of Kadeishvili~\cite{Kadeishvili80}
and Kontsevich--Soibelman~\cite{KontsevichSoibelman00}, in the
form of Loday--Vallette~\cite[\S10.3]{LodayVallette12}, produces
a transferred $\Linf$-algebra structure
$(H^*(\fg),0,\{\ell_n^{\mathrm{tr}}\}_{n\ge 2})$ quasi-isomorphic
to $\fg$, with brackets computed by the Kadeishvili--Merkulov
tree sum
\begin{equation}\label{eq:tree-formula}
\ell_r^{\mathrm{tr}}(a_1,\dots,a_r)
\;=\;
\sum_{T\in\Trees_r}\frac{\pm 1}{|\Aut(T)|}\;
\pi\circ\Phi_T(\iota(a_1),\dots,\iota(a_r)),
\end{equation}
where $\Trees_r$ is the set of planar rooted trees with $r$
leaves and $\Phi_T$ decorates internal vertices with $[-,-]$
and internal edges with $h$.

\subsection{The formality obstruction tower}

\begin{definition}[$\Linf$-formality and obstructions]
\label{def:formality-tower}
The dg Lie algebra $\fg$ is \emph{$\Linf$-formal} if there is an
$\Linf$ quasi-isomorphism $\fg\simeq(H^*(\fg),0,[-,-])$;
equivalently, $\ell_n^{\mathrm{tr}}=0$ in $\Linf$ cohomology for
all $n\ge 3$. The \emph{$\Linf$-formality obstruction tower} is
the sequence
\[
\bigl(o_r(\fg)\bigr)_{r\ge 3},\qquad
o_r(\fg)\;:=\;[\ell_r^{\mathrm{tr}}]\;\in\;
H^*\bigl(\CE^\bullet_{\mathrm{cyc}}(\fg)\bigr).
\]
If $o_r(\fg)=0$, an $\Linf$ homotopy kills
$\ell_r^{\mathrm{tr}}$ and the next obstruction is defined.
\end{definition}

For a pronilpotent dg Lie algebra with Maurer--Cartan element
$\Theta$, evaluating the transferred brackets on $\Theta$
produces a refined tower
$(\ell_r^{\mathrm{tr}}(\Theta,\dots,\Theta))_{r\ge 2}$, which
takes values in the convolution target.  Obstruction cohomology
records the projected class.  The refined target-valued tower is
the one compared with the shadow tower.

\begin{remark}[Cyclic deformation complex]
For $\gmodzero$, the relevant deformation complex is the cyclic
Chevalley complex of the genus-zero convolution dg Lie algebra.
Under the chirally Koszul hypotheses and the chosen cyclic
contraction, this complex maps to the shadow cohomology
$\Ash_{*,0}$ by the projection $\pi$.
\end{remark}

% ================================================================
\section{Base cases: degrees $2$, $3$, $4$}\label{sec:base-cases}
% ================================================================

\begin{proposition}[Shadow--formality identification at low degree]
\label{prop:base-cases}
For every verified chirally Koszul standard algebra in the fixed
genus-zero cyclic transfer setting, the following identities hold:
\begin{enumerate}[label=\textup{(\roman*)}]
\item $\Sh_2(\cA)\;=\;\kappa(\cA)
\;=\;\ell_2^{(0)}(\pi\Theta_{\cA},\pi\Theta_{\cA})$.
\item $\Sh_3(\cA)\;=\;\mathfrak{C}(\cA)
\;=\;-h\bigl(\ell_3^{(0)}(\pi\Theta^{\le 2}_{\cA},
\pi\Theta^{\le 2}_{\cA},\pi\Theta^{\le 2}_{\cA})\bigr)$.
\item $\Sh_4(\cA)\;=\;\mathfrak{Q}(\cA)
\;=\;\bigl[\ell_4^{(0),\mathrm{tr}}(\pi\Theta^{\le 3}_{\cA},
\pi\Theta^{\le 3}_{\cA},\pi\Theta^{\le 3}_{\cA},
\pi\Theta^{\le 3}_{\cA})\bigr]$.
\end{enumerate}
\end{proposition}

\begin{proof}
\emph{Part (i).} The binary bracket on $\gmodzero$ is the
genus-zero convolution Lie bracket; evaluating on
$\Theta_{\cA}$ produces the scalar curvature
$\kappa(\cA)\in\mathbb{Q}(c)$ by
definition~\cite[Theorem~D]{Lorgat26}. The identity
$\Sh_2(\cA)=\kappa(\cA)$ is the degree-$2$ scalar shadow, so the
three quantities agree.

\emph{Part (ii).} The ternary bracket $\ell_3^{(0)}$,
computed by~\eqref{eq:tree-formula}, projects after cyclic
symmetrisation to the three $s$-, $t$-, and $u$-channel stable
four-point trees. On the primary line $[\ell_3^{(0)}]=0$ since
$H^1(\overline{\cM}_{0,4})=H^1(\mathbb{P}^1)=0$; at the chain
level $\ell_3^{(0)}$ is exact and nonzero, and homotopy
transfer gives
$\mathfrak{C}(\cA)=-h(\ell_3^{(0)}(\pi\Theta^{\le 2}_{\cA},
\pi\Theta^{\le 2}_{\cA},\pi\Theta^{\le 2}_{\cA}))$.
The cubic shadow $\mathfrak{C}(\cA)$, computed as the
three-channel graph sum at degree three, matches the cyclic
projection of the transferred tree sum with identical weights and
propagators.

\emph{Part (iii).} The quaternary bracket $\ell_4^{(0)}$ is a
planar HPL tree sum whose cyclic projection is indexed by the
$15$ labelled genus-zero stable trees with five external legs.
Its transferred class
$\ell_4^{(0),\mathrm{tr}}$ lies in $H^*(\overline{\cM}_{0,5})$.
The weight-four graph sum defining $\mathfrak{Q}(\cA)$
enumerates those labelled stable trees with matching weights,
signs, and propagators
\cite[appendix on nonlinear modular shadows]{Lorgat26},
so $\Sh_4(\cA)=\mathfrak{Q}(\cA)=
[\ell_4^{(0),\mathrm{tr}}(\pi\Theta^{\le 3}_{\cA},\dots,
\pi\Theta^{\le 3}_{\cA})]$.
\end{proof}

\begin{remark}[Degree-$2$: $\kappa$ is family-dependent]
\label{rem:kappa-not-universal}
The identity $\Sh_2(\cA)=\kappa(\cA)$ holds at degree two on the
scalar projection for the verified standard families;
on the primary line this is the definition $S_2:=\kappa$.
What is family-dependent is the value of $\kappa$ itself:
$\kappa^{\Vir}=c/2$, $\kappa^{\Heis}=k$,
$\kappa^{\mathrm{KM}}=\dim(\fg)(k+h^\vee)/(2h^\vee)$,
$\kappa^{\cW_N}=c(H_N-1)$ with $H_N=\sum_{j=1}^N 1/j$. The
formula $\kappa=c/2$ is Virasoro-specific.
\end{remark}

% ================================================================
\section{The comparison proof}\label{sec:n6-shadow-formality-proof}
% ================================================================

\begin{theorem}[Genus-zero shadow--formality comparison]
\label{thm:n6-shadow-formality-comparison}
Let $\cA$ be a verified chirally Koszul standard algebra in the
genus-zero cyclic setting, with fixed cyclic contraction
$(\iota,\pi,h)$ of $\gmodzero$ onto $\Ash_{*,0}$. Assume the
cyclic stable-tree/HPL comparison lemma. For every degree
$r\ge 2$,
\begin{equation}\label{eq:n6-shadow-formality-identity}
\Sh_r(\cA)\;=\;\ell_r^{(0),\mathrm{tr}}
\bigl(\pi\Theta^{\le r-1}_{\cA},\ldots,
\pi\Theta^{\le r-1}_{\cA}\bigr)
\quad\text{in }\Ash_{r,0},
\end{equation}
where $\ell_r^{(0),\mathrm{tr}}$ is the transferred genus-zero
$\Linf$ bracket on the minimal model of $\gmodzero$.
For $r=2$ this is the curvature input. For $r\ge3$ it is the
formality obstruction evaluation on the truncated Maurer--Cartan
element.
\end{theorem}

\begin{proof}
Since genus-zero stable graphs are trees, $\gmodzero$ is a
strict dg Lie algebra. The fixed cyclic contraction satisfies
$\pi\iota=\mathrm{id}$, $\mathrm{id}-\iota\pi=[D,h]$. The
homotopy transfer
theorem~\cite{Kadeishvili80,KontsevichSoibelman00,LodayVallette12}
produces a transferred $\Linf$ structure
$(\Ash_{*,0},\{\ell_r^{(0),\mathrm{tr}}\})$, with brackets
computed by~\eqref{eq:tree-formula}.

At $r=2$, both sides equal $\kappa(\cA)$ by
Proposition~\ref{prop:base-cases}(i); the cases $r=3,4$ are
parts (ii), (iii). For $r\ge5$, assume
\eqref{eq:n6-shadow-formality-identity} through degree $r-1$. The
cyclic stable-tree/HPL comparison lemma identifies the graph sum for
$\Sh_r(\cA)$ with the Kadeishvili--Merkulov tree formula for
$\ell_r^{(0),\mathrm{tr}}(\pi\Theta^{\le r-1},\dots,
\pi\Theta^{\le r-1})$: internal edges carry $h$, outputs are
projected by $\pi$, and internal vertices are the lower-degree
transferred brackets. The projected sums agree in $\Ash_{*,0}$.
\end{proof}

\begin{remark}[Genus-zero projection and positive-genus terms]
Theorem~\ref{thm:n6-shadow-formality-comparison} concerns the
genus-zero projection. For scalar uniform-weight families the full
shadow tower has the expansion
\[
\Sh_r(\cA)\;=\;\underbrace{\ell_r^{(0),\mathrm{tr}}
(\pi\Theta^{\le r-1}_{\cA},\dots)}_{\text{genus-zero formality}}
\;+\;\underbrace{\sum_{g\ge 1}\ell_r^{(g),\mathrm{tr}}(\dots)
+\Lambda_P(\dots)}_{\text{genus }\ge 1\text{ corrections}},
\]
where $\Lambda_P$ is the BV loop operator. The first summand is
controlled by classical homotopy transfer. The positive-genus
summands are governed by the quantum $\Linf$ operations and the BV
loop operator. At genus $g\ge 2$ with multi-weight
generators, the scalar formula for the genus corrections
receives a cross-channel term $\delta F_g^{\mathrm{cross}}$;
the genus-zero comparison~\eqref{eq:n6-shadow-formality-identity}
is unchanged.
\end{remark}

% ================================================================
\section{Common Depth}\label{sec:n6-common-depth}
% ================================================================

\begin{definition}[Two genus-zero depth invariants]
For a chiral algebra $\cA$, set
$r_{\max}(\cA):=\sup\{r\ge 2:\Sh_r(\cA)\ne 0\}$ (shadow
termination degree), and
$f_\infty^\Theta(\cA):=\sup\{n\ge 2:
\ell_n^{(0),\mathrm{tr}}(\pi\Theta^{\le n-1}_{\cA},\ldots,
\pi\Theta^{\le n-1}_{\cA})\ne 0\}$
(evaluated $\Linf$ obstruction level on the transferred genus-zero
model).
\end{definition}

\begin{corollary}[Evaluated common depth for the standard cyclic families]
\label{cor:n6-evaluated-depth}
For the verified chirally Koszul standard families in the
genus-zero cyclic setting,
\[
r_{\max}(\cA)=f_\infty^\Theta(\cA).
\]
\end{corollary}

\begin{proof}
The equality is the evaluation identity of
Theorem~\ref{thm:n6-shadow-formality-comparison}, degree by
degree. Operation-level vanishing requires an independent
nondegeneracy lemma for the transferred cyclic primitive component.
\end{proof}

\begin{corollary}[Coarse depth realisation]
\label{cor:n6-coarse-depth-realisation}
Within the constructed standard cyclic families, the four coarse
shadow-depth classes are all realized:
family:
\[
\textbf{G}\ (r_{\max}=2),\quad
\textbf{L}\ (r_{\max}=3),\quad
\textbf{C}\ (r_{\max}=4),\quad
\textbf{M}\ (r_{\max}=\infty).
\]
The class $\textbf{G}$ consists of $\Linf$-formal algebras; the
class $\textbf{M}$ consists of intrinsically non-formal
algebras, in the sense that no finite truncation of their
minimal $\Linf$ model is quasi-isomorphic to a dg Lie algebra.
\end{corollary}

\begin{proof}
Theorem~\ref{thm:n6-shadow-formality-comparison} identifies shadow depth with the
$\Linf$ obstruction level after evaluation on the transferred
Maurer--Cartan element. The examples in Section~\ref{sec:examples}
realize each coarse class.
\end{proof}

% ================================================================
\section{Examples}\label{sec:examples}
% ================================================================

The four standard families realize the four shadow-depth classes.

\begin{example}[Heisenberg; class \textbf{G}, formal]
The Heisenberg vertex algebra $\Heis_k$ at level $k$, with
$J(z)J(w)\sim k/(z-w)^2$, has $\kappa(\Heis_k)=k$. The bar
differential is strictly commutative, so $m_n=0$ and
$\ell_n^{(0),\mathrm{tr}}=0$ for $n\ge 3$. Hence
$r_{\max}(\Heis_k)=2$ and $\Heis_k$ is $\Linf$-formal. Class
$\textbf{G}$.
\end{example}

\begin{example}[Affine Kac--Moody; class \textbf{L}]
At non-critical level $k\ne-h^\vee$, the affine Kac--Moody
vertex algebra $\widehat{\fg}_k$ has
$\kappa(\widehat{\fg}_k)=\dim(\fg)(k+h^\vee)/(2h^\vee)$. The
binary bracket is the Lie bracket on $\fg$; Jacobi kills all
degree-four graph contributions, so $\ell_4^{(0),\mathrm{tr}}=0$,
while $\ell_3^{(0),\mathrm{tr}}\ne 0$ in general. Hence
$r_{\max}(\widehat{\fg}_k)=3$. Class $\textbf{L}$.
\end{example}

\begin{example}[$\beta\gamma_\lambda$; class \textbf{C}, shadow
depth $4$]
The standard rank-one $\beta\gamma_\lambda$ contact family has
\[
S_2=6\lambda^2-6\lambda+1,\qquad
S_3=0,\qquad S_4=-\frac{5}{12},\qquad S_r=0\quad(r\ge5).
\]
Thus $r_{\max}(\beta\gamma_\lambda)=4$ at the generic contact
points. Equivalently, the bracket
$\ell_3^{(0),\mathrm{tr}}$ vanishes and
$\ell_4^{(0),\mathrm{tr}}$ is nonzero. Class $\textbf{C}$.
\end{example}

\begin{example}[Virasoro; class \textbf{M}, obstructed at every
degree]
The Virasoro vertex algebra $\Vir_c$ has $\kappa(\Vir_c)=c/2$.
The transferred brackets on the primary line of
$H^*(\barB(\Vir_c))$ satisfy
$m_k^{\mathrm{tr}}(s^{-1}T,\dots,s^{-1}T)=S_k\cdot e_{2k}$
for $k\ge 2$, with $s^{-1}T$ the desuspension of the conformal
vector to bar degree one and $e_{2k}$ the weight-$2k$ basis vector of
$\Ext^1$. The first shadow coefficients are
\begin{align*}
S_2&=\tfrac{c}{2},&
S_3&=2\ \text{(}c\text{-independent)},&
S_4&=\tfrac{10}{c(5c+22)},\\
S_5&=-\tfrac{48}{c^2(5c+22)},&
S_6&=\tfrac{4(240c+1031)}{c^3(5c+22)^2}.
\end{align*}
The weighted Riccati-metric coefficients $\widehat S^{wt}_6$ and
$\widehat S^{wt}_7$ use a different normalization.
Here $\Ext^1$ denotes the primary bar-cohomology line in the
transferred shadow model. It is not a Hochschild-degree statement;
Theorem H is the separate chiral Hochschild concentration theorem in
degrees $\{0,1,2\}$, with symmetric descent obtained by averaging.
For generic $c$ on the nonsingular Virasoro shadow surface, the
Riccati recurrence and computed coefficients give infinitely many
nonzero $S_k$. Thus $r_{\max}(\Vir_c)=\infty$. The normal-ordered
composites are the source of possible operations, not the
nonvanishing proof. Class $\textbf{M}$.
\end{example}

\begin{remark}[The four standard depths]
For the verified finite-type standard families governed by the
quadratic Riccati reduction, the coarse shadow-depth classes are
$\textbf{G}$, $\textbf{L}$, $\textbf{C}$, and $\textbf{M}$. The
identity $H(t)^2=t^4Q(t)$ with $\deg Q\le2$ gives depths
$2$, $3$, $4$, and $\infty$ for this scalar projection
\cite[theorem on Riccati algebraicity]{Lorgat26}. This is a
coarse classification of the standard cyclic families. Fine
finite-depth towers require separate analysis.
\end{remark}

\begin{thebibliography}{99}

\bibitem{BD04}
A.~Beilinson and V.~Drinfeld.
\emph{Chiral algebras}.
American Mathematical Society Colloquium Publications, 51.
American Mathematical Society, Providence, RI, 2004.

\bibitem{Kadeishvili80}
T.~V. Kadeishvili.
\emph{On the theory of homology of fiber spaces}.
Russian Math. Surveys 35 (1980), no.~3, 231--238.

\bibitem{Kontsevich03}
M.~Kontsevich.
\emph{Deformation quantization of Poisson manifolds}.
Lett. Math. Phys. 66 (2003), no.~3, 157--216.
First circulated as q-alg/9709040 (1997).

\bibitem{KontsevichSoibelman00}
M.~Kontsevich and Y.~Soibelman.
\emph{Homological mirror symmetry and torus fibrations}, in
\emph{Symplectic geometry and mirror symmetry (Seoul, 2000)},
203--263, World Scientific, River Edge, NJ, 2001.

\bibitem{LodayVallette12}
J.-L. Loday and B.~Vallette.
\emph{Algebraic operads}.
Grundlehren der mathematischen Wissenschaften, 346.
Springer, Heidelberg, 2012.

\bibitem{Lorgat26}
R.~Lorgat.
\emph{Modular Koszul Duality, Volume I}.
Monograph, 2026.

\bibitem{Merkulov99}
S.~A. Merkulov.
\emph{Strong homotopy algebras of a K\"ahler manifold}.
Int. Math. Res. Not. 3 (1999), 153--164.

\bibitem{Tamarkin03}
D.~E. Tamarkin.
\emph{Formality of chain operad of little discs}.
Lett. Math. Phys. 66 (2003), no.~1--2, 65--72.

\end{thebibliography}

\end{document}
